\documentclass[10pt,a4paper]{amsart}
\usepackage[margin=19mm]{geometry}
\usepackage{amsmath,amssymb,mathtools}
\usepackage[scr=rsfs,cal=euler]{mathalfa}
\usepackage{xfrac}
\usepackage[unicode,hidelinks,bookmarks=false]{hyperref}
\newcommand{\ie}{i.e.\ }
\newcommand{\cf}{cf.\ }
\newcommand{\resp}{respectively\ }
\newcommand{\Subsection}{\subsection}
\providecommand{\nobreakdash}{\nobreak\hbox{-}}
\setlength{\emergencystretch}{2em}
\multlinegap0pt
\usepackage{bm}
\usepackage{bbm}
\usepackage{dsfont}
\usepackage{xspace}
\usepackage{cancel}
\usepackage{mathtools}
\usepackage{amsthm}
\makeatletter
\@ifundefined{theorem}{\newtheorem{theorem}{Theorem}}{}
\@ifundefined{corollary}{\newtheorem{corollary}[theorem]{Corollary}}{}
\@ifundefined{proposition}{\newtheorem{proposition}[theorem]{Proposition}}{}
\makeatother
\DeclareMathOperator{\Kaz}{Kaz}
\DeclareMathOperator{\Vol}{Vol}
\newcommand{\SchP}{\mathcal{P}}
\newcommand{\SchQ}{\mathcal{Q}}
\title[The formal degree conjecture for groups over local function ]{The formal degree conjecture for groups over local function fields}
\author{Anantha Krishna B}
\date{Source: arXiv:2605.26031v1 (CC BY 4.0)}
\begin{document}
\maketitle

\noindent\emph{Source and licence.} The formal degree conjecture for groups over local function fields. Version arXiv:2605.26031v1 (CC BY 4.0), distributed under the Creative Commons Attribution 4.0 International licence, as recorded in the arXiv licence metadata for this exact version and shown on the abstract page https://arxiv.org/abs/2605.26031v1. Full rights notice: licenses/arxiv-2605.26031-CC-BY-4.0.txt.\par\smallskip

\noindent\textit{Editorial scope.} Counted unit: the Theorem whose source label is \texttt{thm:Formaldeg\_conjecture\_over\_close\_fields} in the article (source0.tex lines 1141--1143, source label \texttt{thm:Formaldeg\_conjecture\_over\_close\_fields}), delivered together with the article's own complete original proof of it (source0.tex lines 1145--1178, one \texttt{proof} environment of 34 lines). No mathematics is written, repaired or re-derived: statement and proof are the article's own text line for line. Required context retained uncounted: Measures\_are\_equal\_close\_field\_type1 (lines 759--762); formal\_degrees\_are\_equal1 (lines 790--795); Thm:Parahoric\_subgroups\_of\_quotients (lines 958--965); param\_over\_close\_fields (lines 1004--1006); L\_functions\_of\_reps\_are\_equal (lines 1075--1077); Eq:dim\_of\_irred\_comp (lines 1133--1135). Every definition, display and label of those lines is retained unchanged. Omitted from this package and not redistributed: 1--758, 763--789, 796--957, 966--1003, 1007--1074, 1078--1132, 1136--1140, 1144--1144, 1179--1426. Nothing inside a retained excerpt is omitted or altered: every retained body is delivered exactly as the article's lines are written, so no omission receipt of any kind accompanies this package. Citations: the retained excerpts carry no citation command at all, so no bibliography entry of the article is transcribed and no reference is redistributed. Reference adaptation: none was needed and none was made; every reference inside the delivered unit resolves inside it, so no unresolved reference or citation is produced. Printed slips kept literal: no printed word, symbol, hypothesis or number of the counted statement or of its proof was altered. Layout and class: the article's own class is \texttt{article} with packages including geometry, inputenc, fontenc, lmodern, textcomp, mathrsfs, authblk, amssymb, latexsym, amsmath, amsthm, mathtools, verbatim, setspace; the wrapper is the lane's pinned \texttt{amsart} editorial wrapper at 10pt on A4, which restates the theorem-like environments the retained text needs and adds the article's own macro definitions that the retained text needs (\texttt{\textbackslash Kaz}, \texttt{\textbackslash SchP}, \texttt{\textbackslash SchQ}, \texttt{\textbackslash Vol}), with extra wrapper packages (bm, bbm, dsfont, xspace, cancel, mathtools). The delivered numbering is the unit's own. Assets: no retained span contains a figure environment, a graphics-inclusion command, a PDF-page inclusion, a table or a TikZ picture; no image file is copied into this package, the package declares no asset dependency and no image file is redistributed with it. Licence: version arXiv:2605.26031v1 of this article is distributed under the Creative Commons Attribution 4.0 International licence, as recorded in the arXiv licence metadata for this exact version and shown on the abstract page https://arxiv.org/abs/2605.26031v1. A full rights notice is delivered at licenses/arxiv-2605.26031-CC-BY-4.0.txt. Characters: every retained excerpt is pure ASCII, so the delivered engine, XeLaTeX with its default Latin Modern text font, reports no missing character.
\begin{corollary} \label{Measures_are_equal_close_field_type1}
    If $\mu_{\psi}(\mathfrak{o}_F) = \mu_{\psi'}(\mathfrak{o}_{F'})$, then
    \[ \int_{P_m} \lvert \omega_G \rvert = \int_{P'_m} \lvert \omega_{G'} \rvert. \]
\end{corollary}

\begin{theorem}\label{formal_degrees_are_equal1}
    Let $\sigma$ and $\sigma'$ correspond via the bijection mentioned above. Then $\sigma \in \Pi_*^2(G,F)_m$ if and only if $\sigma' \in \Pi_*^2(G', F')_m$. In this case,
    \[
    d(\sigma, d\bar{g}) = d(\sigma', d\bar{g}').
    \]
\end{theorem}

\begin{theorem}\label{Thm:Parahoric_subgroups_of_quotients}
    Let $\SchP$, $\SchQ$, and $\mathcal{Z}$ be as above. Then we have:
    \begin{enumerate}
        \item The parahoric subgroups of $G/Z$ attached to vertices are of the form $(\SchP/\mathcal{Z})(\mathfrak{o}) \cong \SchP(\mathfrak{o})/\mathcal{Z}(\mathfrak{o})$.
        \item $\SchP/\mathcal{Z} \cong \SchQ$.
        \item $Q_m \cong P_m Z(F)/Z(F)$.
    \end{enumerate}
\end{theorem}

\begin{proposition}\label{param_over_close_fields}
    The parameter $\varphi$ is discrete if and only if $\varphi'$ is discrete. Furthermore, $\mathcal{S}_{\varphi}^{\sharp} \cong \mathcal{S}_{\varphi'}^{\sharp}$ and $\overline{S_{\varphi}} \cong \overline{S_{\varphi'}}$.
\end{proposition}

\begin{proposition}\label{L_functions_of_reps_are_equal}
    We have $L(\varphi,V,s) = L(\varphi',V',s)$, $\gamma(\varphi,V,s) = \gamma(\varphi',V',s)$, and $\epsilon(\varphi,V,s) = \epsilon(\varphi',V',s)$.
\end{proposition}

\begin{equation}\label{Eq:dim_of_irred_comp}
    \dim(\rho_{\pi}) = \dim(\rho_{\pi'}).
\end{equation}

\begin{theorem}\label{thm:Formaldeg_conjecture_over_close_fields}
    Set $m \ge 1$. There exists $l \ge m$ such that if $F \sim_l F'$, then the formal degree conjecture holds for $(F, G, \varphi, \pi, \psi)$ if and only if it holds for $(F', G', \varphi', \pi', \psi')$.
\end{theorem}

\begin{proof}
Suppose the conductor of $\psi$ is $k$. Then the self-dual Haar measure $\mu_{\psi}$ on $F$ satisfies $\mu_{\psi}(\mathfrak{o}_F) = q^{k/2}$. Because the character $\psi'$ has conductor $k$, we have $\mu_{\psi'}(\mathfrak{o}_{F'}) = q^{k/2}$. Therefore, $\mu_{\psi}(\mathfrak{o}_F) = \mu_{\psi'}(\mathfrak{o}_{F'})$.

By Theorem~\ref{formal_degrees_are_equal1}, there exists $l \ge m$ such that whenever $F \sim_l F'$, the Hecke algebra isomorphism $\Kaz_m \colon \mathcal{H}(G,P_m) \cong \mathcal{H}(G',P_m')$ implies the equality $d(\pi, d\bar{g}) = d(\pi', d\bar{g}')$. There exist $c, c' > 0$ such that $\mu_{Z \backslash G} = c \, d\bar{g}$ and $\mu_{Z' \backslash G'} = c' \, d\bar{g}'$.
We wish to show that
\begin{equation}\label{eq: formal_deg_equal_wrt_gross_measure}
    d(\pi, \mu_{Z \backslash G}) = d(\pi', \mu_{Z' \backslash G'}).
\end{equation}
To show Equation~\ref{eq: formal_deg_equal_wrt_gross_measure}, it is sufficient to prove that $c = c'$.

Recall that $d\bar{g}$ is a measure on $Z(F) \backslash G(F)$ such that $dg = dz \, d\bar{g}$, where $dg$ and $dz$ satisfy $\Vol(P_m; dg) = 1 = \Vol(Z(F) \cap P_m; dz)$. Let $d\nu$ be the measure on $G(F)$ such that $d\nu = dz \, \mu_{Z \backslash G}$, and similarly define $d\nu'$ on $G'(F')$. Then
\[
\Vol(P_m; d\nu) = \Vol(Q_m; \mu_{Z \backslash G}),
\]
where $Q_m = Z(F) \backslash P_m Z(F)$. By Theorem~\ref{Thm:Parahoric_subgroups_of_quotients}, Corollary~\ref{Measures_are_equal_close_field_type1}, and the fact that $a(M_{G/Z}) = a(M_{G'/Z'})$, it follows that
\[
\Vol(Q_m; \mu_{Z \backslash G}) = \Vol(Q_m'; \mu_{(Z\backslash G)'}).
\]
Also, we have
\[
\Vol(Q_m; \mu_{Z \backslash G}) = c \Vol(Q_m; d\bar{g}) = c \Vol(P_m; dg) = c.
\]
Similarly,
\[
\Vol(Q_m'; \mu_{(Z \backslash G)'}) = c'.
\]
Hence, $c = c'$.



 Proposition \ref{param_over_close_fields} and \ref{L_functions_of_reps_are_equal} imply $\# \mathcal S_{\varphi}^{\sharp} = \# \mathcal{S_{\varphi'}^{\sharp}}$ and $\gamma(\varphi) = \gamma(\varphi').$ By our assumption \ref{Eq:dim_of_irred_comp}, it follows that $\dim \rho_{\pi} = \dim \rho_{\pi'}.$ Thus, the theorem follows.
 

\end{proof}

\end{document}
